\documentclass[11pt,a4paper]{article}

\usepackage[UTF8]{ctex}
\usepackage[a4paper,margin=1in]{geometry}
\usepackage{graphicx}
\usepackage{booktabs}
\usepackage{longtable}
\usepackage{float}
\usepackage{hyperref}
\usepackage{array}
\usepackage{caption}
\usepackage{amsmath}
\usepackage{amssymb}

\title{检查点 4 实验报告\\极简 Backbone 搭建、极简 YOLO Loss 与训练跑通}
\author{}
\date{\today}

\begin{document}

\maketitle
\tableofcontents
\newpage

\section{实验目标}

本次检查点 4 的目标是完成一个最小可训练的目标检测链路，重点不在检测精度，而在于把模型、标签编码、损失函数和训练流程真正跑通。具体目标包括：

\begin{enumerate}
  \item 手写一个 10--20 层的极简 CNN backbone。
  \item 手写一个简化版 YOLOv1 风格损失函数，只保留格子分类与 bbox 中心点、宽高回归。
  \item 跑通前向、反向训练流程，并使用 TensorBoard 记录 training loss 曲线。
  \item 在真实样本上可视化当前模型的预测框与真实框对比。
\end{enumerate}

\section{工程结构}

本次实验在 \texttt{/home/lidz/YOLO/miniyolo} 下按真实工程项目方式组织代码，主要目录如下：

\begin{itemize}
  \item \texttt{models/}：模型结构，包含 backbone、detection head 和完整 MiniYOLO。
  \item \texttt{data/}：真实数据读取与标签编码。
  \item \texttt{losses/}：极简 YOLO 损失函数。
  \item \texttt{tools/}：模型检查、训练、测试、可视化与曲线导出脚本。
  \item \texttt{runs/}：TensorBoard 日志目录。
  \item \texttt{outputs/}：模型权重、预测可视化结果与训练曲线图片。
  \item \texttt{logs/}：材料汇总与实验记录。
  \item \texttt{report/}：LaTeX 报告与最终 PDF。
\end{itemize}

\section{模型设计}

\subsection{极简 Backbone}

本次实验设计的 backbone 仅使用最基础的 \texttt{Conv2d + ReLU + MaxPool2d} 组合。输入图像大小统一为 \texttt{224x224}，通过 5 次 \texttt{2x2} 池化，将空间分辨率逐步压缩为 \texttt{7x7}。最终输出特征图大小为：

\[
[B, 512, 7, 7]
\]

之所以选择 \texttt{224} 作为输入尺寸，是因为经过连续 5 次池化后可以稳定得到 \texttt{224 -> 112 -> 56 -> 28 -> 14 -> 7}，便于后续采用 YOLOv1 风格的 \texttt{7x7} 网格预测。

\subsection{Detection Head}

在 backbone 输出特征图 \texttt{[B, 512, 7, 7]} 的基础上，进一步接一个极简 detection head：

\begin{itemize}
  \item 先通过一个 \texttt{3x3} 卷积将通道从 \texttt{512} 压缩到 \texttt{256}
  \item 再通过一个 \texttt{1x1} 卷积得到最终预测维度
\end{itemize}

本实验设定类别数为 \texttt{20}，每个格子同时预测一个 bbox 参数向量 \texttt{(x, y, w, h)}，因此每个格子的输出维度为：

\[
20 + 4 = 24
\]

所以 detection head 最终输出为：

\[
[B, 7, 7, 24]
\]

其中：

\begin{itemize}
  \item 前 20 维表示分类预测
  \item 后 4 维表示 bbox 回归值
\end{itemize}

\subsection{张量形状说明}

利用 \texttt{tools/inspect\_model.py} 对 backbone 与 head 的张量形状进行了逐层检查，结果如表 \ref{tab:backbone-shape} 所示。

\begin{table}[H]
  \centering
  \caption{极简 backbone 的逐层 tensor shape}
  \label{tab:backbone-shape}
  \small
  \begin{tabular}{ccc}
    \toprule
    层编号 & 模块类型 & 输出张量形状 \\
    \midrule
    input & 输入图像 & [2, 3, 224, 224] \\
    layer0 & ConvBlock & [2, 16, 224, 224] \\
    layer1 & ConvBlock & [2, 16, 224, 224] \\
    layer2 & MaxPool2d & [2, 16, 112, 112] \\
    layer3 & ConvBlock & [2, 32, 112, 112] \\
    layer4 & ConvBlock & [2, 32, 112, 112] \\
    layer5 & MaxPool2d & [2, 32, 56, 56] \\
    layer6 & ConvBlock & [2, 64, 56, 56] \\
    layer7 & ConvBlock & [2, 64, 56, 56] \\
    layer8 & MaxPool2d & [2, 64, 28, 28] \\
    layer9 & ConvBlock & [2, 128, 28, 28] \\
    layer10 & ConvBlock & [2, 128, 28, 28] \\
    layer11 & MaxPool2d & [2, 128, 14, 14] \\
    layer12 & ConvBlock & [2, 256, 14, 14] \\
    layer13 & ConvBlock & [2, 256, 14, 14] \\
    layer14 & MaxPool2d & [2, 256, 7, 7] \\
    layer15 & ConvBlock & [2, 512, 7, 7] \\
    layer16 & ConvBlock & [2, 512, 7, 7] \\
    output & backbone 输出 & [2, 512, 7, 7] \\
    head output & 最终预测 & [2, 7, 7, 24] \\
    \bottomrule
  \end{tabular}
\end{table}

进一步将检测输出拆分后，可得到：

\begin{itemize}
  \item 分类预测：\texttt{class\_pred = [B, 7, 7, 20]}
  \item 框回归预测：\texttt{bbox\_pred = [B, 7, 7, 4]}
\end{itemize}

\section{标签编码方式}

\subsection{为什么使用网格坐标}

本实验采用 YOLOv1 风格的网格监督方式。对于每个真实框，不直接用单一的 \texttt{x, y} 监督，而是拆成两部分：

\begin{enumerate}
  \item \texttt{grid\_x, grid\_y}：表示目标中心点落在哪个网格单元
  \item \texttt{cell\_x, cell\_y}：表示目标中心点在当前网格内部的相对偏移
\end{enumerate}

这样可以明确回答两个问题：

\begin{itemize}
  \item 哪个格子负责当前目标
  \item 目标中心点在该格子内部的精确位置
\end{itemize}

\subsection{标签张量设计}

为了与模型预测 \texttt{[B, 7, 7, 24]} 对齐，本实验将监督信息拆成 3 个张量：

\begin{itemize}
  \item \texttt{target\_cls}: \texttt{[7, 7, 20]}
  \item \texttt{target\_box}: \texttt{[7, 7, 4]}
  \item \texttt{object\_mask}: \texttt{[7, 7]}
\end{itemize}

其含义如下：

\begin{itemize}
  \item \texttt{target\_cls}：当前格子的分类 one-hot 监督
  \item \texttt{target\_box}：当前格子的 bbox 回归监督，内容为 \texttt{(cell\_x, cell\_y, w, h)}
  \item \texttt{object\_mask}：该格子是否负责一个目标，取值为 0 或 1
\end{itemize}

\subsection{标签编码过程}

真实框原始格式采用 \texttt{xyxy}，即：

\[
[x_1, y_1, x_2, y_2]
\]

首先转换为：

\[
c_x = \frac{x_1 + x_2}{2}, \quad c_y = \frac{y_1 + y_2}{2}, \quad w = x_2 - x_1, \quad h = y_2 - y_1
\]

然后对 \texttt{c\_x, c\_y, w, h} 做图像大小归一化，再通过 \texttt{grid\_size=7} 映射到目标格子：

\[
grid_x = \lfloor c_x \times 7 \rfloor, \quad grid_y = \lfloor c_y \times 7 \rfloor
\]

格子内部偏移为：

\[
cell_x = c_x \times 7 - grid_x, \quad cell_y = c_y \times 7 - grid_y
\]

最后将：
\begin{itemize}
  \item \texttt{object\_mask[grid\_y, grid\_x] = 1}
  \item \texttt{target\_cls[grid\_y, grid\_x, label] = 1}
  \item \texttt{target\_box[grid\_y, grid\_x] = [cell\_x, cell\_y, w, h]}
\end{itemize}

\section{损失函数设计}

\subsection{损失组成}

本实验的极简 YOLO Loss 只保留最核心的两部分：

\begin{enumerate}
  \item 格子分类损失
  \item bbox 回归损失
\end{enumerate}

不包含：

\begin{itemize}
  \item 多尺度预测
  \item anchor
  \item objectness / confidence 分支
\end{itemize}

\subsection{损失计算方式}

记模型输出为：

\[
pred \in \mathbb{R}^{B \times 7 \times 7 \times 24}
\]

将其拆成：

\[
pred\_cls \in \mathbb{R}^{B \times 7 \times 7 \times 20}, \quad pred\_box \in \mathbb{R}^{B \times 7 \times 7 \times 4}
\]

利用 \texttt{object\_mask} 只在正样本网格上计算损失。分类损失和框损失都采用均方误差（MSE）：

\[
L_{cls} = \frac{\sum ((pred\_cls - target\_cls)^2 \cdot mask)}{N_{pos} \times 20}
\]

\[
L_{box} = \frac{\sum ((pred\_box - target\_box)^2 \cdot mask)}{N_{pos} \times 4}
\]

总损失定义为：

\[
L = \lambda_{cls} \cdot L_{cls} + \lambda_{box} \cdot L_{box}
\]

其中本实验采用：

\[
\lambda_{cls}=1.0,\quad \lambda_{box}=5.0
\]

即：

\[
L = L_{cls} + 5 \cdot L_{box}
\]

这样设置的原因是，bbox 回归通常比分类更难收敛，因此给予更高权重。

\section{数据与训练设置}

\subsection{数据来源}

真实训练数据来自统一后的 manifest：

\begin{itemize}
  \item \texttt{/home/lidz/YOLO/DataSet/Unified/manifests/all\_val.jsonl}
\end{itemize}

本实验为快速验证链路，使用：

\begin{itemize}
  \item \texttt{max\_samples = 16}
\end{itemize}

图片通过 \texttt{DetectionDataset} 读取，并统一 resize 到：

\begin{itemize}
  \item \texttt{224x224}
\end{itemize}

bounding box 也会同步按缩放比例变换。

\subsection{训练超参数}

本次真实训练采用如下设置：

\begin{table}[H]
  \centering
  \caption{训练超参数设置}
  \begin{tabular}{cc}
    \toprule
    项目 & 设置 \\
    \midrule
    数据源 & \texttt{all\_val.jsonl} \\
    最大样本数 & 16 \\
    输入尺寸 & 224 \\
    网格大小 & 7 \\
    类别数 & 20 \\
    batch size & 4 \\
    epochs & 2 \\
    学习率 & 1e-3 \\
    优化器 & Adam \\
    设备 & CUDA \\
    TensorBoard 日志目录 & \texttt{runs/real\_train} \\
    模型权重输出 & \texttt{outputs/miniyolo\_real\_last.pth} \\
    \bottomrule
  \end{tabular}
\end{table}

\section{训练结果与曲线分析}

\subsection{训练结果}

真实数据训练输出的 epoch 级结果如下：

\begin{itemize}
  \item epoch 1:
    \begin{itemize}
      \item total loss = 1.6907
      \item cls loss = 0.0792
      \item box loss = 0.3223
    \end{itemize}
  \item epoch 2:
    \begin{itemize}
      \item total loss = 1.2641
      \item cls loss = 0.0528
      \item box loss = 0.2423
    \end{itemize}
\end{itemize}

可以看到：
\begin{itemize}
  \item 总损失从 1.6907 下降到 1.2641
  \item 分类损失从 0.0792 下降到 0.0528
  \item 框回归损失从 0.3223 下降到 0.2423
\end{itemize}

说明在这组极小规模真实训练中，模型已经表现出初步收敛趋势。

\subsection{Loss 曲线}

\begin{figure}[H]
  \centering
  \includegraphics[width=0.9\textwidth]{../outputs/training_curves/loss_step_curves.png}
  \caption{step 级别 loss 曲线}
\end{figure}

\begin{figure}[H]
  \centering
  \includegraphics[width=0.9\textwidth]{../outputs/training_curves/loss_epoch_curves.png}
  \caption{epoch 级别 loss 曲线}
\end{figure}

\subsection{曲线解释}

本实验导出的 loss 曲线分为两种：

\begin{enumerate}
  \item \textbf{step 曲线}：
  \begin{itemize}
    \item \texttt{loss/total\_step}
    \item \texttt{loss/classification\_step}
    \item \texttt{loss/box\_step}
  \end{itemize}
  这些值是每个 batch 直接计算得到的原始 loss 值，因此波动较大。

  \item \textbf{epoch 曲线}：
  \begin{itemize}
    \item \texttt{loss/total\_epoch}
    \item \texttt{loss/classification\_epoch}
    \item \texttt{loss/box\_epoch}
  \end{itemize}
  这些值是一个 epoch 内所有 batch loss 的平均值，更适合用来观察整体下降趋势。
\end{enumerate}

从曲线可以看出：

\begin{itemize}
  \item step 曲线存在明显抖动，这是因为当前训练样本数很少、batch 也很小
  \item epoch 曲线整体下降，说明模型已经开始学习
  \item 由于总损失定义为 \texttt{cls + 5 * box}，所以总损失的变化主要受 bbox 回归项影响
\end{itemize}

\section{预测框与真实框可视化}

\subsection{可视化方式}

在 \texttt{tools/visualize\_predictions.py} 中，使用训练后权重：

\begin{itemize}
  \item \texttt{outputs/miniyolo\_real\_last.pth}
\end{itemize}

对真实样本进行推理，并将预测框与真实框画在同一张图上。可视化规则如下：

\begin{itemize}
  \item 绿色框：真实框（ground truth）
  \item 绿色文字：\texttt{gt:\{label\}}
  \item 红色框：模型预测框
  \item 红色文字：\texttt{pred:\{class\_id\} \{score\} (\{gy\},\{gx\})}
\end{itemize}

其中：

\begin{itemize}
  \item \texttt{class\_id} 表示预测类别编号
  \item \texttt{score} 表示该预测的分类分数
  \item \texttt{(gy, gx)} 表示该预测来自哪个网格位置
\end{itemize}

为了保证图中至少能看到红色预测框，本次可视化时将分数阈值设为：

\begin{itemize}
  \item \texttt{score\_threshold = 0.0}
\end{itemize}

并设置：

\begin{itemize}
  \item \texttt{top\_k = 3}
\end{itemize}

即每张图最多展示 3 个预测框。

\subsection{可视化结果}

\begin{figure}[H]
  \centering
  \includegraphics[width=0.85\textwidth]{../outputs/prediction_visualizations/voc2012_val_2008_000002.png}
  \caption{样本 \texttt{voc2012\_val\_2008\_000002} 的预测框与真实框对比}
\end{figure}

\begin{figure}[H]
  \centering
  \includegraphics[width=0.85\textwidth]{../outputs/prediction_visualizations/voc2012_val_2008_000003.png}
  \caption{样本 \texttt{voc2012\_val\_2008\_000003} 的预测框与真实框对比}
\end{figure}

\subsection{结果解释}

当前极简模型在可视化结果中的红色预测框较多偏差，和真实框并不吻合，这是当前阶段的正常现象，原因主要包括：

\begin{itemize}
  \item 训练样本数非常少，仅使用了 16 个样本
  \item 训练 epoch 很少，仅训练 2 个 epoch
  \item 模型结构极简，不包含更强的检测设计
  \item 当前损失函数不包含 objectness/confidence 分支
\end{itemize}

因此，本实验阶段的目标不是追求检测精度，而是验证：

\begin{itemize}
  \item 模型可以输出预测框
  \item 预测结果可以成功解码并可视化
  \item 能与真实框放在同一张图上进行对比
\end{itemize}

\section{结论}

本次检查点 4 已完成一个最小可训练目标检测系统的搭建，并验证了从输入图像到检测输出、标签编码、损失计算、训练记录和结果可视化的完整链路。实验结论如下：

\begin{enumerate}
  \item 手写的极简 backbone 能稳定将 \texttt{224x224} 输入压缩到 \texttt{[B, 512, 7, 7]} 特征图。
  \item detection head 能将特征图进一步映射为 \texttt{[B, 7, 7, 24]} 的 YOLO 风格预测张量。
  \item 通过 \texttt{target\_encoder.py}，真实框可以被正确编码为 \texttt{target\_cls}、\texttt{target\_box} 和 \texttt{object\_mask}。
  \item 极简 YOLO Loss 能输出标量 training loss，并支持反向传播。
  \item 在真实数据上训练 2 个 epoch 后，epoch 平均总损失从 1.6907 下降到 1.2641，说明模型已经具备基本学习能力。
  \item 当前预测框质量仍然较差，但已经成功完成与真实框的同图对比，这满足了检查点 4 对“跑通链路与可视化”的核心要求。
\end{enumerate}

\appendix

\section{附录：脚本与工具作用说明}

\begin{longtable}{>{\raggedright\arraybackslash}p{0.32\textwidth}>{\raggedright\arraybackslash}p{0.58\textwidth}}
\toprule
脚本/文件 & 作用 \\
\midrule
\endfirsthead
\toprule
脚本/文件 & 作用 \\
\midrule
\endhead

\path{models/minimal_backbone.py} & 定义极简 CNN backbone，负责把输入图像压缩为 \texttt{7x7} 特征图。 \\
\path{models/detection_head.py} & 定义 detection head，将 backbone 特征图映射为 YOLO 风格预测输出。 \\
\path{models/miniyolo.py} & 将 backbone 与 detection head 封装为完整的 MiniYOLO 模型。 \\
\path{losses/minimal_yolo_loss.py} & 定义极简 YOLO Loss，计算分类损失、bbox 回归损失和总损失。 \\
\path{data/target_encoder.py} & 将真实框和类别编码为 \texttt{target\_cls}、\texttt{target\_box} 和 \texttt{object\_mask}。 \\
\path{data/detection_dataset.py} & 从统一 manifest 读取真实样本、读取图片、缩放 bbox，并生成训练所需监督张量。 \\
\path{tools/inspect_model.py} & 检查 backbone 与 head 的输出 shape，并记录每层 tensor shape。 \\
\path{tools/test_target_encoder.py} & 验证 target encoder 是否能将单个真实框正确写入网格监督张量。 \\
\path{tools/test_minimal_yolo_loss.py} & 验证极简 YOLO Loss 是否能输出标量并支持 \texttt{backward()}。 \\
\path{tools/test_detection_dataset.py} & 验证真实数据版 \texttt{Dataset} 与 \texttt{DataLoader} 是否工作正常。 \\
\path{tools/train_minimal.py} & 使用 toy data 构造最小训练闭环，验证模型、loss 和优化器能协同工作。 \\
\path{tools/train_with_tensorboard.py} & 使用 toy data 训练并写出 TensorBoard loss 日志。 \\
\path{tools/train_real.py} & 使用真实统一数据进行训练，记录 loss 曲线并保存模型权重。 \\
\path{tools/visualize_predictions.py} & 加载训练后权重，在真实样本上画出预测框与真实框对比图。 \\
\path{tools/export_training_curves.py} & 从 TensorBoard 事件文件中导出 step/epoch loss 曲线图片。 \\
\path{TensorBoard} & 浏览训练过程中的 loss 变化曲线。 \\
\path{runs/} & 保存 TensorBoard 事件文件。 \\
\path{outputs/} & 保存训练权重、预测可视化图与 loss 曲线导出图片。 \\
\path{logs/checkpoint4_materials.md} & 汇总 checkpoint 4 当前可直接用于报告的材料。 \\

\bottomrule
\end{longtable}

\end{document}
